\documentclass[11pt]{article}
\usepackage[a4paper,margin=2.5cm]{geometry}
\usepackage{amsmath,amssymb,amsthm}
\usepackage{microtype}
\usepackage[colorlinks=true,linkcolor=blue,urlcolor=blue,citecolor=blue]{hyperref}

\newtheorem{problem}{Problem}

\title{Identifiability and Convexification of a Band-Limited \\
Unassigned Distance-Geometry Problem}
\author{}
\date{}

\begin{document}
\maketitle

Let $X=(x_1,\dots,x_N)\in(\mathbb{R}^3)^N$ carry labels $e_i\in E$ and weights
$b_i>0$. Quotient by rigid motions (translations and rotations, but
\emph{not} reflections):
$\mathcal{X}=(\mathbb{R}^3)^N/E(3)$.
The weighted pair-distance measure and its measurement are
\[
  \rho_X=\sum_{i<j} w_{ij}\,\delta_{\|x_i-x_j\|},\quad w_{ij}=b_i b_j\ge 0,
  \qquad
  F(X)=R_{Q_{\max},\sigma}\,\rho_X\in Y,
\]
where $R_{Q_{\max},\sigma}$ is a known compact smoothing operator: a
band-limit $Q_{\max}$ (real-space resolution $\Delta r\approx\pi/Q_{\max}$)
followed by convolution of width $\sigma>0$. Its \emph{information capacity}
\[
  m_{\mathrm{eff}}(Q_{\max},\sigma)
  =\dim\operatorname{span}\{R_{Q_{\max},\sigma}\delta_t:t\in[0,R_{\max}]\}
  \approx\lceil R_{\max}/\Delta r\rceil
\]
is finite, with $m_{\mathrm{eff}}\to\infty$ as $Q_{\max}\to\infty$
(exact-distance limit) and $m_{\mathrm{eff}}\to 0$ as $Q_{\max}\to 0$, and
$\operatorname{rank}(DF_X)\le\min(d_C,m_{\mathrm{eff}})$.
A \emph{prior} is a closed $C\subset\mathcal{X}$, a smooth manifold of
dimension $d_C$ realizable in internal coordinates (a graph of fixed edge
lengths and angles with $d_C$ free dihedrals). $F$ is an instance of the
unassigned distance geometry problem \cite{Billinge2018}; in the motivating
application $C$ encodes chemical constraints, which play no formal role below.

\paragraph{Known facts.}
\begin{enumerate}
\item For $R=\mathrm{id}$, injectivity of $\rho$ up to rigid motion is the
global-rigidity problem; in $3$D generic global rigidity requires redundant
rigidity and stressability and is subtler than edge-counting, and special
configurations are generically \emph{not} globally rigid (flip ambiguities
persist).
\item In $1$D (turnpike/beltway) with unassigned noisy distances, projected
gradient descent with a spectral initializer converges locally linearly to
the global optimizer under explicit conditions \cite{HuangDokmanic2021}; for
integer sets, sparsity yields polynomial-time recovery up to shift and
reversal w.h.p.\ \cite{JaganathanHassibi2012}.
\item A configuration and its mirror image have identical $F(X)$, so
identifiability holds at best modulo $\mathbb{Z}_2$ (fiber $K\ge2$ for chiral
configurations).
\end{enumerate}

\paragraph{Open problem.}
For $y=F(X_\star)$ and $\mathcal{L}_y(X)=\|F(X)-y\|_Y^2|_C$:
\begin{itemize}
\item[\textbf{(P1)}] \emph{Identifiability:} $F|_C$ injective modulo
reflection (fibers of size $1$, or $2$ = the mirror pair).
\item[\textbf{(P2)}] \emph{Local identifiability:} the Fisher operator
$I_X=(DF_X)^{*}(DF_X)$ has full rank on $T_XC$ for all $X\in C$.
\item[\textbf{(P3)}] \emph{$\varepsilon$-convexity:} $\mathcal{L}_y|_C$ is
geodesically $\varepsilon$-convex on the component of $X_\star$ with a unique
stationary point (no spurious local minima).
\end{itemize}

\begin{problem}
Characterize the priors $C$, as a function of $m_{\mathrm{eff}}$, for which
$\text{(P1)}\!\Leftrightarrow\!\text{(P2)}\!\Leftrightarrow\!\text{(P3)}$
hold for generic $X_\star\in C$, and find the \emph{minimal} such prior. In
particular:
\begin{enumerate}
\item[\textbf{(a)}] Prove or disprove a sharp threshold at
$m_{\mathrm{eff}}=d_C$ (modulo reflection), with counterexamples at
$m_{\mathrm{eff}}=d_C-1$.
\item[\textbf{(b)}] Find the maximal admissible $d_C$ (least external
information) and the trade-off curve $d_C^{\,*}(Q_{\max})$.
\item[\textbf{(c)}] Show the threshold coincides with vanishing of
$I_{\mathrm{deficit}}=h(X\mid C)-I(F(X);X\mid C)$ (up to $\log 2$).
\end{enumerate}
\end{problem}

This is generic-global-rigidity theory for band-limited, weighted, soft data
on a constraint manifold --- the regime open at the intersection of
\cite{HuangDokmanic2021,JaganathanHassibi2012,Billinge2018}.

\begin{thebibliography}{99}
\bibitem{HuangDokmanic2021} S.~Huang, I.~Dokmani\'c, \emph{Reconstructing
Point Sets from Distance Distributions}, IEEE Trans.\ Signal Process.\
\textbf{69} (2021), 1811--1827. arXiv:1804.02465.
\bibitem{JaganathanHassibi2012} K.~Jaganathan, B.~Hassibi, \emph{Reconstruction
of Integers from Pairwise Distances}, arXiv:1212.2386 (2012).
\bibitem{Billinge2018} S.~J.~L.~Billinge, P.~M.~Duxbury, D.~S.~Gon\c{c}alves,
C.~Lavor, A.~Mucherino, \emph{Recent results on assigned and unassigned
distance geometry with applications to protein molecules and nanostructures},
Ann.\ Oper.\ Res.\ \textbf{271}(1) (2018), 161--203. DOI:10.1007/s10479-018-2989-6.
\bibitem{Billinge2004} S.~J.~L.~Billinge, M.~G.~Kanatzidis, \emph{Beyond
crystallography: the study of disorder, nanocrystallinity and
crystallographically challenged materials with pair distribution functions},
Chem.\ Commun.\ (2004), 749--760. DOI:10.1039/b311073n.
\end{thebibliography}
\end{document}
